% Lösungen Einheit 1 von 3 – Koordinatensystem (zusammenbau v0.1)
\einheitenkopf{Einheit 1 von 3 · Koordinatensystem}

\erg{10}{a) Pfeil 4 Kästchen nach rechts, daran Pfeil 1 Kästchen nach oben}
%% TODO Lösung der Erklärzeile 11g) fehlt in der Bank
\erg{11}{a) zuerst nach rechts \quad b) $A(4|6)$: 4 nach rechts, 6 nach oben \quad c) $B(1|3)$: 1 nach rechts, 3 nach oben \quad d) $C(6|1)$: 6 nach rechts, 1 nach oben \quad e) $D(5|5)$: 5 nach rechts, 5 nach oben \quad f) $E(2|7)$: 2 nach rechts, 7 nach oben \quad g) \ldots \quad h) $F(6|4)$ \quad i) $B(-4|1)$ und $C(2|-3)$ eingetragen (negative Zahl: nach links bzw. nach unten) \quad j) $ 2.$ Quadrant (links oben) \quad k) $K(0|-4)$ liegt auf der Hochachse (y-Achse), denn die erste Koordinate ist 0. \quad l) Rechteck \quad m) $D(1|5)$ \quad n) $D(6|0)$}
\erg{12}{a) $2$}
\erg{13}{a) $Q(6|1)$}
\erg{14}{a) $7 - 1 = 6$ Kästchen}
\erg{15}{a) Sie hat die Koordinaten vertauscht; richtig ist $A(2|6)$: erst 2 nach rechts, dann 6 nach oben. \quad b) Die zweite Koordinate sagt, wie weit man nach oben oder unten geht; auf der Rechtsachse geht man gar nicht nach oben, also ist sie 0. \quad c) $A(1|4)$, $B(5|6)$, $C(6|-2)$ \quad d) $8 \cdot 5 = 40$\,m}
